\section{Method}
\label{sec:method}

\subsection{Data}
\label{sec:data}

Table~\ref{tab:data} summarizes the Criteo Uplift v2.1 benchmark
\cite{diemert2021}. It merges several advertisers' incrementality tests; in
each, a random share of users is held out from ad targeting. Version 2.1
re-sampled every test to a common treatment ratio to remove a leak in v1.
Rows were sub-sampled non-uniformly for privacy, so all rates describe the
benchmark sample, not real-world ad incrementality.

\begin{table}[H]
\centering
\caption{The Criteo Uplift v2.1 dataset as used here.}
\label{tab:data}
\small
\begin{tabular}{l p{9.5cm}}
\toprule
\textbf{Item} & \textbf{Value} \\
\midrule
Users & 13,979,592; split 60/20/20 into 8.39M train, 2.80M validation,
        2.80M test, stratified on treatment $\times$ visit $\times$ conversion \\
Features & 12 anonymized: 4 continuous, 8 hashed categoricals with 60--3,743 levels \\
\texttt{treatment} & randomized ad eligibility; 85.0\,\% treated \\
\texttt{exposure} & shown $\ge 1$ ad; 3.60\,\% of treated users, 0 of control
                    (post-randomization, self-selected) \\
\texttt{visit} / \texttt{conversion} & 4.70\,\% / 0.29\,\% overall \\
\bottomrule
\end{tabular}
\end{table}

\subsection{Pipeline}
\label{sec:pipeline}

Figure~\ref{fig:pipeline} shows the stages. Each reads a single
configuration file (seeds, splits, sample sizes, model settings, business
assumptions) and writes metrics, tables, and figures to disk; the narrative
notebooks only read those outputs.

\begin{figure}[H]
\centering
\resizebox{\linewidth}{!}{\begin{tikzpicture}[
    >=Stealth,
    box/.style={
      rectangle,
      rounded corners=3pt,
      draw=black!60,
      fill=blue!4,
      minimum height=1.45cm,
      minimum width=2.4cm,
      align=center,
      inner sep=4pt,
      text width=2.2cm,
      font=\small
    },
    arrow/.style={->, thick, draw=black!60, shorten <=1pt, shorten >=1pt},
    node distance=0.35cm
  ]

  \node[box] (data)    {\textbf{Data}\\[2pt]\footnotesize 14.0M users,\\ 60/20/20 split};
  \node[box, right=of data]    (eda)     {\textbf{EDA + tests}\\[2pt]\footnotesize rates, balance,\\ C2ST};
  \node[box, right=of eda]     (pred)    {\textbf{Predict}\\[2pt]\footnotesize LR, RF, LightGBM,\\ calibration};
  \node[box, right=of pred]    (ate)     {\textbf{ATE}\\[2pt]\footnotesize DiM, Lin, IPW,\\ AIPW, IV};
  \node[box, below=0.6cm of data]   (cate)    {\textbf{CATE / uplift}\\[2pt]\footnotesize S/T/X/DR,\\ causal forest};
  \node[box, right=of cate]    (policy)  {\textbf{Policy}\\[2pt]\footnotesize budgeted\\ targeting};
  \node[box, right=of policy]  (robust)  {\textbf{Robustness}\\[2pt]\footnotesize seeds, halves,\\ learning curve};
  \node[box, right=of robust]  (scale)   {\textbf{Scaling}\\[2pt]\footnotesize time, memory,\\ rows};

  \draw[arrow] (data) -- (eda);
  \draw[arrow] (eda)  -- (pred);
  \draw[arrow] (pred) -- (ate);
  \draw[arrow] (ate.south) -- ++(0,-0.3) -| ($(data.south)!0.5!(cate.north)$) -- (cate.north);
  \draw[arrow] (cate) -- (policy);
  \draw[arrow] (policy) -- (robust);
  \draw[arrow] (robust) -- (scale);

\end{tikzpicture}
}
\caption{Analysis pipeline. Every stage uses the same stratified
train/validation/test split.}
\label{fig:pipeline}
\end{figure}

\subsection{Protocol}
\label{sec:protocol}

Features, calibrators, thresholds, hyperparameters, budget choices, and
learner selection use only the train and validation splits. The test split
is used once, for the reported numbers. Categorical levels seen fewer than
1,000 times in the training split are pooled into an ``other'' level, and
the feature builder never sees a label.

Uncertainty comes from three sources: paired Poisson bootstraps over test
users (uplift curves and policies), influence functions (ATE estimators),
and refits with new seeds and disjoint training halves (robustness). Unless
stated, intervals are 95\,\%.

\subsection{Predictive models}
\label{sec:predictive}

We fit $P(\text{conversion}\mid X)$ with L2 logistic regression
(interpretable baseline), a random forest on a 2M-row sample, and LightGBM
on all 8.39M training rows with early stopping on validation
\cite{ke2017}. Model selection uses five-fold cross-validated log loss, a
proper scoring rule that is far less noisy than PR-AUC with about 0.3\,\%
positives. Platt and isotonic calibrators are fit on validation and
evaluated on test.

\subsection{Average treatment effects}
\label{sec:ate-method}

Write $T$ for assignment, $Y$ for an outcome, and $X$ for features. The
estimators are the difference in means, Lin regression adjustment
\cite{lin2013} with HC2 standard errors, IPW (Horvitz--Thompson and
H{\'a}jek) with design, logistic, and LightGBM propensities, and the doubly
robust AIPW estimator~\cite{robins1994},
\begin{equation}
  \hat\tau_{\mathrm{AIPW}} = \frac{1}{n}\sum_{i=1}^{n}
  \Big[\hat\mu_1(X_i)-\hat\mu_0(X_i)
  + \frac{T_i\,(Y_i-\hat\mu_1(X_i))}{\hat e(X_i)}
  - \frac{(1-T_i)\,(Y_i-\hat\mu_0(X_i))}{1-\hat e(X_i)}\Big],
\end{equation}
with cross-fitted LightGBM outcome models $\hat\mu_t$ and propensity
$\hat e$; AIPW is our primary estimator. For the effect of actually seeing
an ad, we use assignment $T$ as an instrument for exposure $D$, with a Wald
estimator and a covariate-adjusted variant~\cite{imbens1994}. It estimates
the effect on the exposed under the exclusion restriction that assignment
affects outcomes only through ads shown.

\subsection{Heterogeneous effects and uplift}
\label{sec:cate-method}

We implement S-, T-, X-, and DR-learners on LightGBM base learners, and fit
an honest causal forest with EconML \cite{econml} on a 1M-row training
sample. The transformed-outcome method is included as a class-variable
baseline \cite{athey2015}. Every learner is compared with a
\emph{response model} $P(Y\mid X)$ that ignores treatment. Rankings are
scored by the normalized Qini coefficient, the area between the Qini curve
and random targeting divided by the area of a perfect ranking, with
propensity-weighted bootstrap intervals. We also report the BLP and GATES
diagnostics~\cite{chernozhukov2018gates}.

\subsection{Targeting policies}
\label{sec:policy-method}

Given a budget $b$ (share of users treated), each policy treats the top $b$
by its score: random, highest conversion probability, highest estimated
uplift, and highest expected value, which treats every user with
$\hat\tau(X)\,V - c > 0$ for value $V$ and cost $c$. Policy value is
estimated off-policy with propensity weights on the test split. For
illustration, $V=\$50$ per conversion and $c=\$0.05$ per treated user; these
are assumptions, not measurements.
